\chapter{Offers, Compensation and Non-Competes}\label{ch:iv:offers-compensation-and-non-competes}

Two offers arrive in the same week. One pays a higher base salary and a guaranteed first-year bonus, with a
sign-on payment repayable if the candidate leaves within two years and a twelve-month non-compete. The other
pays less cash now, defers a large share of the bonus over three years, and has no non-compete at all. Which is
worth more depends on numbers the letters do not print: how likely the candidate is to leave, when, and what
the non-compete costs in the year it bites. This chapter reads an offer as a set of cash flows and conditions,
values it, and negotiates it. It explains how the clauses work and where the rules are; it gives no legal
advice, and a candidate with a contract to sign should have it read by a lawyer where they live.

\section{Reading an offer letter}

\begin{definition}[Offer letter]\label{def:iv:offers-compensation-and-non-competes:letter}
An \emph{offer letter}\index{offer letter} is the firm's written proposal of employment: the role, the start
date, the pay and its conditions, the notice period and the restrictive covenants, and the conditions on which
the offer depends (references, background checks, regulatory approval). It is followed by, or incorporates,
the employment contract, whose terms govern.
\end{definition}

The pay vocabulary is One Quant Book~17's (chapter~13: base salary, total compensation, sign-on bonus,
guaranteed bonus, restricted stock units, forfeiture and good-leaver provisions) and the pay-policy and
contract vocabulary is One Quant Book~16's (chapter~10: bonus pool, discretionary bonus, deferred compensation,
vesting schedule, malus and clawback; chapter~11: non-compete clause, garden leave, non-solicitation clause,
confidentiality agreement, invention assignment clause).

\begin{method}[Reading an offer, line by line]\label{met:iv:offers-compensation-and-non-competes:read}
\begin{enumerate}
\item \emph{Fixed pay}: base salary, when it is reviewed, in what currency.
\item \emph{Variable pay}: guaranteed or discretionary; if discretionary, how the pool is set and whether a
formula applies; when it is paid.
\item \emph{Deferral}: what share of variable pay is deferred, into what (cash, shares, fund units), vesting
when, forfeited in which leaving cases (the good-leaver provision defines the exceptions), subject to malus
and clawback.
\item \emph{One-off payments}: sign-on bonus, buy-out of forfeited deferrals from the current employer, and
the conditions for repaying them.
\item \emph{Time}: start date, notice period on both sides, probation.
\item \emph{Covenants}: non-compete (length, scope, whether paid), garden leave, non-solicitation, notice of
competing offers; and whether the firm can waive them.
\item \emph{Conditions}: what must happen before the offer becomes binding.
\end{enumerate}
\end{method}

\section{Valuing an offer}

An offer is a set of cash flows whose amounts and timing depend on whether and when the candidate leaves.
Guaranteed cash is worth what it says; deferred cash is worth what it says only if the candidate stays until it
vests; a sign-on payment is worth nothing if it is repaid; and a non-compete costs the pay of the months it
keeps the candidate out of work, unless those months are paid.

\begin{proposition}[Value of an offer with a leaving probability]\label{prop:iv:offers-compensation-and-non-competes:value}
Let $V_{\mathrm{stay}}$ be the pay received over a horizon if the candidate stays, and $V_{\mathrm{leave}}$ the
pay if the candidate leaves at a given date, counting forfeited deferrals as lost, repaid sign-on payments as
lost, and months kept out of work by a non-compete as unpaid. If the candidate leaves with probability
$\pi_{\mathrm{leave}}$, the offer is worth $V_{\mathrm{offer}} = V_{\mathrm{stay}} - \pi_{\mathrm{leave}}(V_{\mathrm{stay}}
- V_{\mathrm{leave}})$ in expectation. Two offers are equally valuable at
$\pi^\ast = (V^A_{\mathrm{stay}} - V^B_{\mathrm{stay}})/\big((V^A_{\mathrm{stay}} - V^A_{\mathrm{leave}}) -
(V^B_{\mathrm{stay}} - V^B_{\mathrm{leave}})\big)$.
\end{proposition}

\begin{example}[The hook's two offers]\label{ex:iv:offers-compensation-and-non-competes:two}
In thousands, undiscounted and before tax, over three years. Offer A: base 230, a guaranteed first-year bonus
of 120, an expected bonus of 90 in later years, a sign-on of 80 repayable in full on leaving within two years,
a twelve-month unpaid non-compete, no deferral. Offer B: base 170, an expected bonus of 150 a year of which
40\% is deferred and vests in equal thirds over the following three years, no sign-on, no non-compete. If the
candidate stays three years, A pays $80 + 3 \times 230 + 120 + 2 \times 90 = 1\,070$ and B pays $3 \times 170 +
3 \times 150 = 960$. If the candidate leaves after one year for a job paying 320 a year, A pays $230 + 120 =
350$ for the first year, the sign-on is repaid and the non-compete costs the second year: $350 + 320 = 670$;
B pays $170 + 0.6 \times 150 = 260$ for the first year, forfeits the year's 60 of deferral, and nothing else:
$260 + 2 \times 320 = 900$. A loses 400 on leaving and B loses 60, so B is worth more once the probability
of leaving exceeds $110/340 \approx 0.32$ (\cref{fig:iv:offers-compensation-and-non-competes:value}).
\end{example}

\begin{omfigure}
\begin{tikzpicture}
\begin{axis}[width=10cm,height=5.2cm,xlabel={\small probability of leaving after one year},
  ylabel={\small expected pay (thousands)},xmin=0,xmax=1,ymin=600,ymax=1100,
  tick label style={font=\footnotesize},grid=major,grid style={gray!20},table/col sep=comma,
  legend cell align=left,legend style={font=\scriptsize,at={(0.5,-0.32)},anchor=north,
  /tikz/every even column/.append style={column sep=8pt}},legend columns=2]
\addplot[omDef,thick] table[x=p,y=A]{figdata/interviews/07-offers-compensation-and-non-competes/offers.csv};
\addplot[omThm,thick,dashed] table[x=p,y=B]{figdata/interviews/07-offers-compensation-and-non-competes/offers.csv};
\draw[gray,dotted] (axis cs:0.3235,600) -- (axis cs:0.3235,1100);
\legend{offer A (cash now; sign-on; non-compete),offer B (deferral; no non-compete)}
\end{axis}
\end{tikzpicture}

\omcaption{Expected three-year pay of the two offers of
\cref{ex:iv:offers-compensation-and-non-competes:two} against the probability of leaving after one year
(thousands, undiscounted, before tax). The dotted line marks the break-even probability, $11/34 \approx
0.32$. Data: \texttt{fig\_iv\_offer.py}.}\label{fig:iv:offers-compensation-and-non-competes:value}
\end{omfigure}

The same model makes the other trade-offs explicit: discounting shortens the value of deferrals, tax changes
the cash but rarely the ranking, and a higher next salary makes a non-compete more expensive. One Quant
Book~17's \texttt{firm.payoffer} builds a fuller version with tax and vesting schedules; the chapter's own
checks are in \texttt{iv\_offer.py}.

\section{Negotiating}

Negotiation is the subject of One Quant Book~16, chapter~23: each side has a best alternative to a negotiated
agreement, a reservation value it will not go beyond, and between the two reservation values lies the zone of
possible agreement. For a candidate, the best alternative is the next best offer or the current job; the
firm's reservation value is set by its pay bands, what it pays similar hires and what losing the candidate
costs.

\begin{definition}[Exploding offer, counteroffer]\label{def:iv:offers-compensation-and-non-competes:exploding}
An \emph{exploding offer}\index{exploding offer} is an offer that lapses if not accepted within a short
deadline, set to prevent the candidate from completing other processes. A \emph{counteroffer}\index{counteroffer}
is an improved offer made in response to another: by the hiring firm when the candidate asks for more, or by
the current employer when the candidate resigns.
\end{definition}

\begin{method}[Negotiating an offer]\label{met:iv:offers-compensation-and-non-competes:negotiate}
\begin{enumerate}
\item Know your alternative and value it with \cref{prop:iv:offers-compensation-and-non-competes:value}.
\item Learn the firm's range: published pay ranges where the law requires them
(\cref{dat:iv:offers-compensation-and-non-competes:law}), the recruiter, peers.
\item Ask for specific items with reasons: a buy-out of deferrals you will forfeit, a sign-on to cover a
notice-period gap, a guarantee for a first year whose bonus cannot yet reflect your work, a shorter or paid
non-compete.
\item Negotiate once, in writing, with a single consolidated request; accept or decline promptly after.
\item Never invent an offer; never accept an offer you do not intend to honour.
\end{enumerate}
\end{method}

An \omterm{def:iv:offers-compensation-and-non-competes:exploding}{exploding offer} is a negotiating position, not a law of nature: the answer is to ask for a date that lets
the other processes finish (\cref{ch:iv:the-process-by-firm-type}). A \omterm{def:iv:offers-compensation-and-non-competes:exploding}{counteroffer} from the current employer
deserves the same valuation as any offer, plus one question the numbers do not answer: whether the reasons
for leaving are removed by money.

\begin{dated}{2026-09}{Pay transparency and salary history}\label{dat:iv:offers-compensation-and-non-competes:law}
\textbf{European Union}: Directive (EU) 2023/970, article 5, gives applicants the right to receive the initial
pay or its range for the position before the interview or otherwise, and forbids employers to ask applicants
about their pay history; member states had to transpose it by 7 June 2026. \textbf{New York City}: since 1
November 2022, employers with four or more employees advertising a job performed in the city must state the
minimum and maximum salary they in good faith believe they would pay. \textbf{California}: Labor Code section
432.3 forbids employers to seek or rely on an applicant's salary history, requires the pay scale to be given
to an applicant on reasonable request, and requires employers with fifteen or more employees to include it in
job postings.
\end{dated}

\section{Restrictive covenants}

A non-compete, a garden leave clause and a long notice period all delay the day the candidate can work for a
competitor, and they differ in who pays for the delay: garden leave and notice are paid, an unpaid non-compete
is not. Their enforceability differs between jurisdictions and changes over time; the mechanisms and the
litigation record are in One Quant Book~16, chapter~11, and the career consequences in One Quant Book~17,
chapter~28.

\begin{dated}{2026-09}{Where the rules on non-competes stand}\label{dat:iv:offers-compensation-and-non-competes:noncompete}
\textbf{United States, federal}: the Federal Trade Commission's Non-Compete Clause Rule of 2024 is not in effect:
a district court set it aside on 20 August 2024 (Ryan, LLC v.\ FTC), and in September 2025 the Commission voted
to dismiss its appeals and accede to the vacatur. \textbf{California}: Business and Professions Code section
16600.5, in force since 1 January 2024, makes a non-compete that is void under California law unenforceable
regardless of where and when it was signed, and forbids employers to enter into or attempt to enforce one;
section 16600.1 required employers to notify affected employees by 14 February 2024 that such clauses were
void. \textbf{United Kingdom}: the government published a working paper on options for reforming non-compete
clauses on 26 November 2025 (a statutory limit on length, a ban, or a salary threshold); it closed for
responses on 18 February 2026.
\end{dated}

\section{Question bank}

\begin{interviewq}[$\star$ \iqroles{trader, researcher, developer} \iqfirm{any}]\label{iq:iv:offers-compensation-and-non-competes:1}
Which facts must an \omterm{def:iv:offers-compensation-and-non-competes:letter}{offer letter} give you before you can compare it with another offer? List them in the
order you would check them.
\end{interviewq}

\begin{interviewq}[$\star$ \iqroles{trader} \iqfirm{proprietary firm}]\label{iq:iv:offers-compensation-and-non-competes:2}
Your sign-on bonus of 60\,000 is repayable pro rata if you leave within 24 months. You leave after 9 months.
How much do you repay?
\end{interviewq}

\begin{interviewq}[$\star$ \iqroles{researcher} \iqfirm{multi-manager fund}]\label{iq:iv:offers-compensation-and-non-competes:3}
Your bonus of 200\,000 is 40\% deferred, vesting in three equal instalments 12, 24 and 36 months after the
bonus date, and forfeited if you resign. What cash do you receive at the bonus date, what are the instalments,
and what do you forfeit if you resign 18 months after the bonus date?
\end{interviewq}

\begin{interviewq}[$\star\star$ \iqroles{trader} \iqfirm{market maker}]\label{iq:iv:offers-compensation-and-non-competes:4}
Using the two offers of \cref{ex:iv:offers-compensation-and-non-competes:two}, which do you prefer if you
think there is a one-in-five chance you will want to leave after a year? And a one-in-two chance? At what
probability are you indifferent?
\end{interviewq}

\begin{interviewq}[$\star\star$ \iqroles{developer} \iqfirm{proprietary firm}]\label{iq:iv:offers-compensation-and-non-competes:5}
A firm makes you an offer on Tuesday that expires on Thursday. You are in the \omterm{def:iv:the-final-round-and-trading-games:final}{final round} elsewhere next
week. What do you say to the first firm, and what do you do if it refuses to move?
\end{interviewq}

\begin{interviewq}[$\star\star$ \iqroles{researcher, trader} \iqfirm{multi-manager fund}]\label{iq:iv:offers-compensation-and-non-competes:6}
One offer has a six-month garden leave paid at your base of 250 a year; another has a twelve-month unpaid
non-compete. Your next job would pay 400 a year. What does each clause cost you if you leave, in thousands?
Which would you negotiate on, and how?
\end{interviewq}

\begin{interviewq}[$\star\star$ \iqroles{trader, researcher, developer} \iqfirm{bank}]\label{iq:iv:offers-compensation-and-non-competes:7}
The recruiter asks: ``What number do you have in mind, and what are you paid now?'' How do you answer both
questions?
\end{interviewq}

\begin{interviewq}[$\star\star\star$ \iqroles{trader, bank} \iqfirm{bank}]\label{iq:iv:offers-compensation-and-non-competes:8}
You resign to take an offer; your current employer counters by raising your base by 20\% and promising a
larger bonus. How do you decide? Name the quantities you would compare and the one you cannot compute.
\end{interviewq}

\begin{interviewq}[$\star\star\star$ \iqroles{researcher, trader} \iqfirm{multi-manager fund}]\label{iq:iv:offers-compensation-and-non-competes:9}
Your best alternative is worth 260. You believe the firm's walk-away value is uniformly distributed between
250 and 350, and that you can make a single take-it-or-leave-it request: if it is at or below the firm's
walk-away value it is accepted; otherwise you take your alternative. What do you ask for, and what is it worth
to you in expectation? What does the model leave out?
\end{interviewq}

\begin{interviewq}[$\star\star\star$ \iqroles{risk, bank} \iqfirm{bank}]\label{iq:iv:offers-compensation-and-non-competes:10}
A clause reads: ``For twelve months after termination the Employee shall not be engaged in any business that
competes with any business of the Group in any country in which the Group operates.'' What does it restrict,
what questions do you ask before signing, and why might its enforceability depend on where you work?
\end{interviewq}

\begin{omsources}
\item Directive (EU) 2023/970 on pay transparency, articles 5 and 34.
\item California Labor Code section 432.3; California Business and Professions Code sections 16600.1 and
16600.5.
\item New York City Commission on Human Rights, \emph{Salary Transparency in Job Advertisements} (fact sheet).
\item Federal Trade Commission, Non-Compete Clause Rule page (status after Ryan, LLC v.\ FTC).
\item Department for Business and Trade, \emph{Reform of non-compete clauses in employment contracts: working
paper}, 26 November 2025.
\item One Quant Book~16, chapters 10, 11 and~23; One Quant Book~17, chapters 13 and~28.
\end{omsources}
